%-----------------------------------------------------------------------------
% Kamil Krzywoń — CV
% ML / AI Engineer & Data Scientist
% Single-column, ATS-friendly, modern LaTeX template
%
% Kompilacja: npm run cv  (z katalogu site/)
% Wynik: public/cv.pdf — plik JEST commitowany, obraz Dockera nie ma TeX-a.
%-----------------------------------------------------------------------------
\documentclass[10pt, a4paper]{article}

% ----- Packages -----
\usepackage[utf8]{inputenc}
\usepackage[T1]{fontenc}
\usepackage[english]{babel}
\usepackage[margin=1.2cm]{geometry}
\usepackage{enumitem}
\usepackage{titlesec}
\usepackage{xcolor}
\usepackage{hyperref}
\usepackage{fontawesome5}
\usepackage{tabularx}
\usepackage{ragged2e}
\usepackage{microtype}
\usepackage{lmodern}
\usepackage[default,scale=0.95]{opensans}
\renewcommand{\familydefault}{\sfdefault}

% ----- Colors -----
\definecolor{accent}{HTML}{0E4D7A}     % deep blue — name & section rules
\definecolor{subtle}{HTML}{4A4A4A}     % body text
\definecolor{muted}{HTML}{7A7A7A}      % secondary info (dates, locations)
\color{subtle}

% ----- Hyperlink style -----
\hypersetup{
    colorlinks=true,
    urlcolor=accent,
    linkcolor=accent,
    pdftitle={Kamil Krzywoń — CV},
    pdfauthor={Kamil Krzywoń}
}

% ----- Page setup -----
\pagestyle{empty}
\setlength{\parindent}{0pt}
\setlength{\parskip}{0pt}

% ----- Section formatting -----
\titleformat{\section}
    {\color{accent}\large\bfseries}
    {}{0pt}{\MakeUppercase}[\vspace{1pt}\color{accent}\titlerule\vspace{2pt}]
\titlespacing*{\section}{0pt}{6pt}{4pt}

% ----- Helper commands -----
\newcommand{\entry}[4]{%
    \noindent\begin{tabularx}{\textwidth}{@{}X r@{}}
        \textbf{#1} & \textcolor{muted}{#2} \\
        \textit{#3} & \textcolor{muted}{\textit{#4}} \\
    \end{tabularx}\vspace{2pt}
}

\newcommand{\project}[3]{%
    \noindent\textbf{#1}\,\textcolor{muted}{\,---\, #2}\\[1pt]
    #3\vspace{3pt}
}

\newcommand{\skillrow}[2]{%
    \noindent\makebox[3.5cm][l]{\textbf{#1}}#2\\[2pt]
}

\setlist[itemize]{leftmargin=1.2em, topsep=2pt, itemsep=1pt, parsep=0pt}

%-----------------------------------------------------------------------------
\begin{document}

% ============================== HEADER ==============================
\begin{center}
    {\Huge\bfseries\color{accent} Kamil Krzywoń}\\[4pt]
    {\large\color{muted} ML / AI Engineer \,|\, Data Scientist}\\[6pt]
    \small
    \faEnvelope[regular]\ \href{mailto:kamil.krzywon24@gmail.com}{kamil.krzywon24@gmail.com}
    \;\;\textbar\;\;
    \faPhone\ +48 667 865 281
    \;\;\textbar\;\;
    \faMapMarker*\ Cieszyn, Poland
    \\[2pt]
    \faGithub\ \href{https://github.com/Kamilox007}{https://github.com/Kamilox007}
    \;\;\textbar\;\;
    \faLinkedin\ \href{https://www.linkedin.com/in/kamil-krzywo\%C5\%84-349354314/}{linkedin.com/in/kamil-krzywoń}
\end{center}

\vspace{4pt}

% ============================== SUMMARY ==============================
\section*{Profile}
Final-year Computer Science student at Poznań University of Technology specialising in
Artificial Intelligence. Hands-on experience building machine learning models in Python
with PyTorch, TensorFlow, and scikit-learn, with a growing focus on NLP and large language
models.

% ============================== EDUCATION ==============================
\section*{Education}

\entry{Poznań University of Technology}{Oct 2021 -- Sep 2026 (expected)}
      {M.Sc. \& B.Sc., Computer Science --- specialisation: Artificial Intelligence}{Poznań, Poland}
\vspace{2pt}
\begin{itemize}
    \item \textbf{Master's thesis (in progress):} \emph{Robust value-based method for group decision support.}
    \item \textbf{Relevant coursework:} Machine Learning, Deep Learning, Natural Language Processing,
          Probability \& Statistics, Linear Algebra, Optimization, Algorithms \& Data Structures, Databases.
\end{itemize}

% ============================== EXPERIENCE ==============================
\section*{Experience}

\entry{bit4mation Polska Sp. z o.o.}{Aug 2024}
      {Software Engineering Intern}{Poznań, Poland}
% DO UZUPEŁNIENIA — nie znam szczegółów tego stażu. Dwa punkty wystarczą:
% czego dotyczył projekt, w jakim stacku, co konkretnie zrobiłeś.
% \begin{itemize}
%     \item ...
%     \item ...
% \end{itemize}
\vspace{3pt}

\entry{Echo Edukacji Sp. z o.o.}{Oct 2024 -- Jun 2025}
      {Math and Physics Tutor}{Poznań, Poland}
% DO UZUPEŁNIENIA — np. liczba uczniów, poziom (matura / studia), wyniki.
\vspace{3pt}

\entry{Private Tutoring (Self-Employed)}{Jun 2025 -- Present}
      {Math and Physics Tutor}{Poznań, Poland}
\begin{itemize}
    \item Teach mathematics and physics at secondary level, including final-exam preparation.
    \item Built and operate the scheduling and billing platform used to run the practice
          (see \emph{Projects})
\end{itemize}

% ============================== SKILLS ==============================
\section*{Technical Skills}

\skillrow{Programming}{Python (advanced), SQL}
\skillrow{ML \& Deep Learning}{PyTorch, TensorFlow, scikit-learn, Hugging Face Transformers}
\skillrow{Data \& Visualisation}{pandas, NumPy, matplotlib, seaborn}
\skillrow{Backend \& Web}{FastAPI, SQLAlchemy, Alembic, REST APIs, React}
\skillrow{MLOps \& Tools}{Git, Docker, MLflow, pytest, Linux, \LaTeX}
\skillrow{Databases}{PostgreSQL, MongoDB, SQLite}
\skillrow{Cloud}{AWS (basic)}
\skillrow{Focus areas}{Classical ML, NLP, Large Language Models, Reinforcement Learning}

% ============================== PROJECTS ==============================
\vspace{-12pt}
\section*{Projects}

\project{Tutoring Management Platform --- production system}
        {Python, FastAPI, SQLAlchemy, React, Docker, Caddy}
        {Full-stack application for lesson scheduling, billing, and student accounts,
         in daily production use. JWT auth with role-based access control, account lockout,
         per-IP rate limiting. Deployed via Docker
         Compose behind Caddy on a Linux VPS, with Alembic migrations, continuous database
         replication to object storage (Litestream), and a regression test suite.}

\project{6 nimmt! — Reinforcement Learning Card-Game Agent}
        {Python, PyTorch}
        {Built an AI agent that learns to play the card game \emph{6 nimmt!} using
         Deep Q-Networks (DQN). Designed the game environment, state representation, and
         reward shaping; trained and evaluated multiple bot variants in self-play. University project.}

\project{Machine Learning Mini-Projects}
        {Python, scikit-learn, pandas, PyTorch}
        {Collection of smaller end-to-end ML projects covering classification, regression,
         and exploratory data analysis on tabular and text data --- including data cleaning,
         feature engineering, model selection, and evaluation with cross-validation.}

% ============================== LANGUAGES ==============================
\section*{Languages}
\skillrow{Polish}{Native}
\skillrow{English}{B2 / C1 (upper-intermediate to advanced)}
% ============================== INTERESTS ==============================
\vspace{-12pt}

\section*{Interests}

ML enthusiast \,$\bullet$\, Mathematics \,$\bullet$\, Physics \,$\bullet$\,
Video games \,$\bullet$\, Board games

\end{document}